%% Solutions for Chapter 10: Pushdown Automata

\chapter{Pushdown Automata --- Solutions}

\begin{enumerate}[{10.}1. ]

%%------------------------------------------------------------
\item \textbf{Show by induction: $\DBAR(s_0,x)=t \Leftrightarrow \langle s_0,x,\cent\rangle\Dvdash\langle t,\lambda,\cent\rangle$.}

\textbf{Claim} ($P(n)$): For all $x \in \Sigma^n$, $\DBAR(s_0,x)=t$ iff $\langle s_0,x,\cent\rangle\Dvdash\langle t,\lambda,\cent\rangle$.

\textbf{Base} ($n=0$): $\DBAR(s_0,\lambda)=s_0$ and $\langle s_0,\lambda,\cent\rangle\Dvdash\langle s_0,\lambda,\cent\rangle$ (zero moves). Both sides give $t=s_0$. \checkmark

\textbf{Step}: Assume $P(n)$. For $x=ya$ ($|y|=n$, $a\in\Sigma$):

($\Rightarrow$): $\DBAR(s_0,ya)=\delta(\DBAR(s_0,y),a)=t$. Let $q=\DBAR(s_0,y)$. By IH, $\langle s_0,y,\cent\rangle\Dvdash\langle q,\lambda,\cent\rangle$. By Theorem 10.1, $\langle q,a,\cent\rangle\vdash\langle t,\lambda,\cent\rangle$ (single step reading $a$). So $\langle s_0,ya,\cent\rangle\Dvdash\langle t,\lambda,\cent\rangle$.

($\Leftarrow$): If $\langle s_0,ya,\cent\rangle\Dvdash\langle t,\lambda,\cent\rangle$, then the computation passes through some state $q$ after reading $y$: $\langle s_0,y,\cent\rangle\Dvdash\langle q,\lambda,\cent\rangle$. By IH, $\DBAR(s_0,y)=q$. The last step reads $a$ from state $q$ to reach $t$, so $\delta(q,a)=t$, hence $\DBAR(s_0,ya)=t$. \qed

%%------------------------------------------------------------
\item \textbf{Define a one-state DPDA $P'_1$ accepting $\{\pmb{a}^n\pmb{b}^n \mid n\geq 0\}$ via final state.}

A one-state DPDA accepting via final state always accepts $\lambda$: the unique state must be both the initial state and the only final state, so the machine is in a final state at the very start with empty input. Hence the best a one-state final-state DPDA can achieve is $\{\pmb{a}^n\pmb{b}^n \mid n\geq 0\}$ (which includes $\lambda$).

Let $P'_1 = \langle\{\pmb{a},\pmb{b}\},\{Z,A\},\{t\},t,\delta,Z,\{t\}\rangle$ where:
\begin{align*}
\delta(t,\pmb{a},Z) &= \langle t, AZ\rangle \quad\text{(push $A$ on reading first \pmb{a})} \\
\delta(t,\pmb{a},A) &= \langle t, AA\rangle \quad\text{(push another $A$ for each subsequent \pmb{a})} \\
\delta(t,\pmb{b},A) &= \langle t, \lambda\rangle \quad\text{(pop $A$ on reading \pmb{b})}
\end{align*}
The machine pushes one $A$ for each $\pmb{a}$ read, then pops one $A$ for each $\pmb{b}$ read. It accepts when the input is exhausted in state $t$: this occurs precisely when the $\pmb{a}$s and $\pmb{b}$s are balanced (stack returned to $Z$) or when the input is empty. Note: $\delta(t,\pmb{b},Z)$ is undefined, so a string with more $\pmb{b}$s than $\pmb{a}$s causes the machine to get stuck (reject). Hence $\Lambda(P'_1) = \{\pmb{a}^n\pmb{b}^n \mid n\geq 0\}$.

%%------------------------------------------------------------
\item \textbf{PDA $P'_2$: proofs and language.}

$P'_2 = \langle\{\pmb{a},\pmb{b}\},\{S,C\},\{t\},t,\delta,S,\{t\}\rangle$ with $\delta(t,\pmb{a},S) = \{\langle t,SC\rangle,\langle t,C\rangle\}$, $\delta(t,\pmb{b},C) = \{\langle t,\lambda\rangle\}$, others empty.

\begin{enumerate}
\item \textbf{Inductive proof}: $(\forall i\in\NN)$: if $\langle t,\pmb{a}^i,S\rangle\Dvdash\langle t,\lambda,\alpha\rangle$ then $\alpha = SC^i$ or $\alpha = C^i$.

\textbf{Base} ($i=0$): $\langle t,\lambda,S\rangle\Dvdash\langle t,\lambda,S\rangle$; here $\alpha=S=SC^0$. \checkmark

\textbf{Step}: Assume the result for $i$. Reading $\pmb{a}^{i+1}$ from $S$: first move on $(\pmb{a},S)$ either pushes $SC$ or $C$.

Case 1: Push $SC$. Then $\langle t,\pmb{a}^{i+1},S\rangle\vdash\langle t,\pmb{a}^i,SC\rangle$. Now process $\pmb{a}^i$ from $S$: by IH, $\langle t,\pmb{a}^i,S\rangle\Dvdash\langle t,\lambda,\beta\rangle$ where $\beta\in\{SC^i,C^i\}$. So $\langle t,\pmb{a}^i,SC\rangle\Dvdash\langle t,\lambda,\beta C\rangle\in\{SC^{i+1},C^{i+1}\}$.

Case 2: Push $C$. Then $\langle t,\pmb{a}^{i+1},S\rangle\vdash\langle t,\pmb{a}^i,C\rangle$. Processing $\pmb{a}^i$ from $C$: there are no moves (neither $\delta(t,\pmb{a},C)$ nor $\delta(t,\lambda,C)$ are defined for reading $\pmb{a}$), so this branch is a dead end unless $i=0$.

So $\alpha\in\{SC^{i+1},C^{i+1}\}$. \qed

\item \textbf{Inductive proof}: $(\forall i\in\NN)$: if $\langle t,x,C^i\rangle\Dvdash\langle t,\lambda,\beta\rangle$ then $x=\pmb{b}^i$.

\textbf{Base} ($i=0$): $\langle t,x,\lambda\rangle$ can only have $x=\lambda$ (stack empty, no moves possible). So $x=\lambda=\pmb{b}^0$. \checkmark

\textbf{Step}: Assume for $i$. From $\langle t,x,C^{i+1}\rangle$: the top of the stack is $C$. The only move consuming $C$ reads $\pmb{b}$: $\delta(t,\pmb{b},C)=\langle t,\lambda\rangle$. So $x$ must start with $\pmb{b}$: $x=\pmb{b}x'$ and $\langle t,\pmb{b}x',C^{i+1}\rangle\vdash\langle t,x',C^i\rangle$. By IH, $x'=\pmb{b}^i$. Hence $x=\pmb{b}^{i+1}$. \qed

\item $P'_2$ accepts by empty stack (the final-state component $\{t\}$ is not used for acceptance here; a string is accepted iff the machine empties its stack while in state $t$ with empty input). From (a), processing $\pmb{a}^n$ from start symbol $S$ can yield stack $C^n$ (by choosing $\delta(t,\pmb{a},S)=\langle t,SC\rangle$ at each of the first $n-1$ steps and then $\langle t,C\rangle$ at the last, so the final stack is $C^n$). From (b), the stack $C^n$ is emptied by reading exactly $\pmb{b}^n$. Thus:
\[
L(P'_2) = \{\pmb{a}^n\pmb{b}^n \mid n\geq 1\}.
\]
For $n=0$: the initial configuration $\langle t,\lambda,S\rangle$ has $S$ still on the stack with no moves available, so the stack is never emptied and $\lambda\notin L(P'_2)$. Hence $L(P'_2) = \{\pmb{a}^n\pmb{b}^n \mid n\geq 1\}$.
\end{enumerate}

%%------------------------------------------------------------
\item \textbf{$L=\{\pmb{a}^i\pmb{b}^j\pmb{c}^k \mid i,j,k\in\NN,\ i+j=k\}$.}

\begin{enumerate}
\item \textbf{PDA via final state}: Use a single counting symbol $A$: push one $A$ for each $\pmb{a}$ or $\pmb{b}$, then pop one $A$ for each $\pmb{c}$. A machine is
\[
P=\langle\{\pmb{a},\pmb{b},\pmb{c}\},\{Z,A\},\{q_0,q_1,q_2,q_f\},q_0,\delta,Z,\{q_0,q_f\}\rangle,
\]
where the nonempty transitions are
\begin{align*}
\delta(q_0,\pmb{a},Z)&=\langle q_0,AZ\rangle, &
\delta(q_0,\pmb{a},A)&=\langle q_0,AA\rangle, \\
\delta(q_0,\pmb{b},Z)&=\langle q_1,AZ\rangle, &
\delta(q_0,\pmb{b},A)&=\langle q_1,AA\rangle, \\
\delta(q_0,\pmb{c},A)&=\langle q_2,\lambda\rangle, &
\delta(q_1,\pmb{b},A)&=\langle q_1,AA\rangle, \\
\delta(q_1,\pmb{c},A)&=\langle q_2,\lambda\rangle, &
\delta(q_2,\pmb{c},A)&=\langle q_2,\lambda\rangle, \\
\delta(q_2,\lambda,Z)&=\langle q_f,Z\rangle. &&
\end{align*}
State $q_0$ handles the $\pmb{a}$-block, $q_1$ the $\pmb{b}$-block, and $q_2$ the $\pmb{c}$-block. The string $\lambda$ is accepted because $q_0$ is final; otherwise, acceptance occurs exactly when all pushed $A$'s have been matched by $\pmb{c}$'s.

\item \textbf{PDA via empty stack}: Use the same three reading states and replace the last move by $\delta(q_2,\lambda,Z)=\langle q_2,\lambda\rangle$. To include $\lambda$, also allow $\delta(q_0,\lambda,Z)=\langle q_0,\lambda\rangle$.

\item \textbf{Counting automaton}: Yes. In fact the PDA in part (a) already uses only one genuine counting symbol $A$, so it is a counting automaton.

\item \textbf{DPDA}: Yes. At each stage there is only one possible action: push during the $\pmb{a}$- and $\pmb{b}$-blocks, then pop during the $\pmb{c}$-block. Thus the machine above can be made deterministic.

\item \textbf{Grammar from PDA} (using Definition 10.7): Applying Definition 10.7 to part (a) and deleting useless variables yields the simpler equivalent grammar
\[
S \rightarrow \pmb{a}S\pmb{c}\mid T,\qquad
T \rightarrow \pmb{b}T\pmb{c}\mid \lambda.
\]
This generates exactly the strings $\pmb{a}^i\pmb{b}^j\pmb{c}^{i+j}$.
\end{enumerate}

%%------------------------------------------------------------
\item \textbf{$L=\{x\in\{\pmb{a},\pmb{b},\pmb{c}\}^* \mid |x|_{\pmb{a}}+|x|_{\pmb{b}}=|x|_{\pmb{c}}\}$.}

\begin{enumerate}
\item \textbf{PDA via final state}: Use two stack symbols $A$ (surplus of $\pmb{a}$/$\pmb{b}$) and $B$ (surplus of $\pmb{c}$) to handle symbols in any order.

$P$: single reading state $q$, final state $f$, bottom marker $Z$:
\begin{align*}
\delta(q,\pmb{a},Z)&=\langle q,AZ\rangle, & \delta(q,\pmb{a},A)&=\langle q,AA\rangle, & \delta(q,\pmb{a},B)&=\langle q,\lambda\rangle,\\
\delta(q,\pmb{b},Z)&=\langle q,AZ\rangle, & \delta(q,\pmb{b},A)&=\langle q,AA\rangle, & \delta(q,\pmb{b},B)&=\langle q,\lambda\rangle,\\
\delta(q,\pmb{c},Z)&=\langle q,BZ\rangle, & \delta(q,\pmb{c},B)&=\langle q,BB\rangle, & \delta(q,\pmb{c},A)&=\langle q,\lambda\rangle.
\end{align*}
When $\pmb{a}$ or $\pmb{b}$ is read with $B$ on top, the pending $\pmb{c}$-surplus is cancelled, and vice versa.
Accept by final state: $\delta(q,\lambda,Z)=\langle f,Z\rangle$ (fires when input exhausted and stack balanced).

\item \textbf{PDA via empty stack}: Pop $Z$ as well. Add $\delta(q,\lambda,Z)=\langle q,\lambda\rangle$ (the stack becomes empty when balanced).

\item \textbf{Counting automaton}: Yes. One counter is enough to store the current difference between $|x|_{\pmb{a}}+|x|_{\pmb{b}}$ and $|x|_{\pmb{c}}$, while the finite control records which side currently has the surplus.

\item \textbf{DPDA}: Yes. The two-symbol construction in part~(a) is deterministic: for each $(q, \sigma, X)$ triple there is at most one transition, since every input symbol and every stack top symbol determine a unique action. The only potential conflict is the $\lambda$-transition $\delta(q,\lambda,Z)=\langle f,Z\rangle$, which must not coexist with other $q$-transitions on $Z$. This is resolved by moving the $\lambda$-acceptance to a separate final state reached only when input is exhausted; in practice, the construction is deterministic for all reachable configurations. Hence a DPDA for $L$ exists.

\item \textbf{Grammar from PDA}: Applying Definition 10.7 to part (a) and simplifying gives an equivalent grammar such as
\[
S\rightarrow \pmb{a}S\pmb{c}\mid \pmb{c}S\pmb{a}\mid \pmb{b}S\pmb{c}\mid
\pmb{c}S\pmb{b}\mid SS\mid \lambda.
\]
Each production preserves the invariant $|x|_{\pmb{a}}+|x|_{\pmb{b}}=|x|_{\pmb{c}}$, and every nonempty balanced word can be split at some point into one balanced word or into a balanced word bracketed by one of the four possible outer pairs.
\end{enumerate}

%%------------------------------------------------------------
\item \textbf{Closure of \PSIG\ and \ASIG\ under inverse homomorphism.}

\begin{enumerate}
\item \textbf{\PSIG\ is closed under inverse homomorphism}: Given a context-free language $L$ over $\Sigma$ and a homomorphism $h: \Delta^* \to \Sigma^*$, we want $h^{-1}(L) = \{w \in \Delta^* \mid h(w) \in L\}$ to be a CFL.

\textbf{Proof}: Construct a PDA for $h^{-1}(L)$ from the PDA $P$ for $L$. The new PDA $P'$ uses a buffer: for each input symbol $a\in\Delta$, it replaces $a$ with $h(a)$ in a buffer and then simulates $P$ on $h(a)$ before reading the next input symbol. The new PDA processes $h(a)$ letter-by-letter from the buffer. Since $h$ maps to a finite string, this is a finite process per input symbol. Thus $P'$ accepts exactly $h^{-1}(L)$, which is a CFL.

\item \textbf{\ASIG\ (DCFL) is closed under inverse homomorphism}: Given a DCFL $L$ and a homomorphism $h$, $h^{-1}(L)$ is also a DCFL.

\textbf{Proof}: The same buffering construction applied to a DPDA for $L$ yields a DPDA for $h^{-1}(L)$, since the DPDA is deterministic and the buffering of $h(a)$ (a fixed string per input symbol $a$) preserves determinism. \qed
\end{enumerate}

%%------------------------------------------------------------
\item \textbf{A finite language not recognized by any one-state PDA via final state.}

\textbf{Example}: $L = \{\pmb{ab}\}$.

\textbf{Proof}: Every one-state PDA $P$ with state set $\{t\}$, start state $t$, and final states $\{t\}$ trivially accepts $\lambda$: the initial configuration has $P$ already in the (unique) final state with empty input, so $\lambda$ is accepted regardless of the stack content. Since $\lambda \notin L = \{\pmb{ab}\}$, no one-state PDA via final state can recognize $L$.

More generally, any finite language that does not contain $\lambda$ is not recognizable by any one-state PDA via final state, since such a machine always accepts $\lambda$. \qed

%%------------------------------------------------------------
\item \textbf{$L=\{\pmb{a}^n\pmb{b}^n\pmb{c}^m\pmb{d}^m \mid n,m\in\NN\}$.}

\begin{enumerate}
\item \textbf{PDA via final state}: Push $A$ on $\pmb{a}$; pop $A$ on $\pmb{b}$; push $B$ on $\pmb{c}$; pop $B$ on $\pmb{d}$; accept when balanced.
States: $\{q_0,q_1,q_2,q_f\}$. Initial state $q_0$, final state $q_f$. Stack alphabet $\{Z,A,B\}$.
$\delta(q_0,\pmb{a},Z)=\langle q_0,AZ\rangle$, $\delta(q_0,\pmb{a},A)=\langle q_0,AA\rangle$,
$\delta(q_0,\pmb{b},A)=\langle q_1,\lambda\rangle$, $\delta(q_1,\pmb{b},A)=\langle q_1,\lambda\rangle$,
After reading $\pmb{b}^n$, the machine is in state $q_1$ with $Z$ on top of the stack. Then:
$\delta(q_1,\pmb{c},Z)=\langle q_2,BZ\rangle$, $\delta(q_2,\pmb{c},B)=\langle q_2,BB\rangle$,
$\delta(q_2,\pmb{d},B)=\langle q_2,\lambda\rangle$, $\delta(q_2,\lambda,Z)=\langle q_f,Z\rangle$.
Also $\delta(q_0,\lambda,Z)=\langle q_1,Z\rangle$ to handle $n=0$, and $\delta(q_1,\lambda,Z)=\langle q_2,Z\rangle$ for $m=0$ (with appropriate final states).

\item \textbf{PDA via empty stack}: Similar; pop $Z$ at the end.

\item \textbf{DPDA}: Yes. The DPDA deterministically pushes $A$s, then pops $A$s, then pushes $B$s, then pops $B$s. Each step is deterministic.

\item \textbf{Counting automaton}: Yes. A single stack symbol suffices: use it first to count the $\pmb{a}$'s against the $\pmb{b}$'s and, after the stack returns to $Z$, reuse the same symbol to count the $\pmb{c}$'s against the $\pmb{d}$'s.

\item \textbf{Grammar from PDA} (Definition 10.7): After applying Definition 10.7 to part (b) and deleting useless variables, we get the equivalent grammar
\[
S\rightarrow AB,\qquad
A\rightarrow \pmb{a}A\pmb{b}\mid \lambda,\qquad
B\rightarrow \pmb{c}B\pmb{d}\mid \lambda.
\]
\end{enumerate}

%%------------------------------------------------------------
\item \textbf{Prove by induction (Theorem 10.2): $\langle s,x,S\rangle\Dvdash\langle s,\lambda,\beta\rangle \Leftrightarrow S\DRightarrow x\beta$ as a leftmost derivation.}

\textbf{Claim} ($P(n)$): For all $n\geq 0$ and strings $x,\beta$: $\langle s,x,S\rangle\Dvdash\langle s,\lambda,\beta\rangle$ in $n$ moves iff $S\DRightarrow x\beta$ in $n$ leftmost-derivation steps.

\textbf{Base} ($n=0$): $\langle s,\lambda,S\rangle\Dvdash\langle s,\lambda,S\rangle$ (zero moves) iff $x=\lambda, \beta=S$, and $S\DRightarrow^0 \lambda\cdot S$ (zero steps). \checkmark

\textbf{Step}: Assume $P(n)$ for all shorter computations. A computation of $n+1$ steps:
$\langle s,ay,S\rangle\vdash\langle s,y,\gamma\rangle\Dvdash\langle s,\lambda,\beta\rangle$ (first move reads $a$, pushes $\gamma$ for top of $S$).
By the construction of the PDA (Theorem 10.2), this first step corresponds to applying a grammar production $S \to a\gamma$ (GNF production, starting with terminal $a$). By $P(n)$, $\langle s,y,\gamma\rangle\Dvdash\langle s,\lambda,\beta\rangle$ iff $\gamma\DRightarrow y\beta$ as leftmost derivation.
So the full computation corresponds to $S \Rightarrow a\gamma \DRightarrow ay\beta = x\beta$. \qed

%%------------------------------------------------------------
\item \textbf{Grammar for regular expressions: convert to GNF, find corresponding PDA.}

Grammar: $R \to \pmb{a}|\pmb{b}|\pmb{c}|\pmb{\epsilon}|\pmb{\emptyset}|(R\pmb{\cdot}R)|(R\pmb{\cup}R)|R^{\pmb{*}}$.

\begin{enumerate}
\item \textbf{GNF conversion}: Since pure Greibach normal form only requires the first symbol on the right to be terminal, the only problem is the left-recursive rule $R\rightarrow R\pmb{^*}$. Introduce a new nonterminal $Y$ generating one or more copies of $\pmb{^*}$:
\[
Y\rightarrow \pmb{^*}\mid \pmb{^*}Y.
\]
Then an equivalent PGNF grammar is
\begin{align*}
R \rightarrow\;& \pmb{a}\mid \pmb{b}\mid \pmb{c}\mid \pmb{\epsilon}\mid \pmb{\emptyset}\mid
\pmb{(}R\pmb{\cdot}R\pmb{)}\mid \pmb{(}R\pmb{\cup}R\pmb{)} \\
\mid\;& \pmb{a}Y\mid \pmb{b}Y\mid \pmb{c}Y\mid \pmb{\epsilon}Y\mid \pmb{\emptyset}Y\mid
\pmb{(}R\pmb{\cdot}R\pmb{)}Y\mid \pmb{(}R\pmb{\cup}R\pmb{)}Y .
\end{align*}
This simply replaces the old rule $R\rightarrow R\pmb{^*}$ by attaching a nonempty string of stars through $Y$.

\item \textbf{PDA from GNF grammar} (Definition 10.6): Let
\[
P=\langle\Sigma,\Gamma,\{s\},s,\delta,R,\emptyset\rangle,
\]
where
\[
\Sigma=\{\pmb{a},\pmb{b},\pmb{c},\pmb{(},\pmb{)},\pmb{\epsilon},\pmb{\emptyset},
\pmb{\cup},\pmb{\cdot},\pmb{^*}\}
\]
and $\Gamma=\Sigma\cup\{R,Y\}$. The nonempty transitions are exactly those dictated by Definition 10.6:
\begin{align*}
\delta(s,\pmb{a},R)&=\{\langle s,\lambda\rangle,\langle s,Y\rangle\}, \\
\delta(s,\pmb{b},R)&=\{\langle s,\lambda\rangle,\langle s,Y\rangle\}, \\
\delta(s,\pmb{c},R)&=\{\langle s,\lambda\rangle,\langle s,Y\rangle\}, \\
\delta(s,\pmb{\epsilon},R)&=\{\langle s,\lambda\rangle,\langle s,Y\rangle\}, \\
\delta(s,\pmb{\emptyset},R)&=\{\langle s,\lambda\rangle,\langle s,Y\rangle\}, \\
\delta(s,\pmb{(},R)&=\{\langle s,R\pmb{\cdot}R\pmb{)}\rangle,\langle s,R\pmb{\cup}R\pmb{)}\rangle,
\langle s,R\pmb{\cdot}R\pmb{)}Y\rangle,\langle s,R\pmb{\cup}R\pmb{)}Y\rangle\}, \\
\delta(s,\pmb{^*},Y)&=\{\langle s,\lambda\rangle,\langle s,Y\rangle\},
\end{align*}
and for each terminal $\tau\in\Sigma$,
\[
\delta(s,\tau,\tau)=\{\langle s,\lambda\rangle\}.
\]
This PDA accepts via empty stack.

\item \textbf{PDA accepting via final state} (Theorem 10.4): Apply Theorem 10.4 to the machine in part (b). Using fresh bottom markers $U$ and $V$, we obtain
\[
P_f=\langle\Sigma,\Gamma\cup\{U,V\},\{s,f\},s,\delta_f,V,\{f\}\rangle,
\]
where $\delta_f$ agrees with $\delta$ on all old stack symbols, and the additional nonempty transitions are
\[
\delta_f(s,\lambda,V)=\{\langle s,RU\rangle\},\qquad
\delta_f(s,\lambda,U)=\{\langle f,U\rangle\}.
\]
There are no other moves involving $U$ or $V$, and no moves out of $f$. By Theorem 10.4, $L(P_f)=\Lambda(P)$.
\end{enumerate}

%%------------------------------------------------------------
\item \textbf{$L=\{\pmb{a}^i\pmb{b}^j\pmb{c}^j\pmb{d}^i \mid i,j\in\NN\}$.}

\begin{enumerate}
\item \textbf{PDA via final state}: Push $A$ on $\pmb{a}$; push $B$ on $\pmb{b}$; pop $B$ on $\pmb{c}$; pop $A$ on $\pmb{d}$; accept when stack returns to $Z$.
States $\{q_0,q_1,q_2,q_3,q_f\}$, alphabet $\{Z,A,B\}$:
$\delta(q_0,\pmb{a},Z)=\langle q_0,AZ\rangle$, $\delta(q_0,\pmb{a},A)=\langle q_0,AA\rangle$,
$\delta(q_0,\pmb{b},A)=\langle q_1,BA\rangle$, $\delta(q_0,\pmb{b},Z)=\langle q_1,BZ\rangle$ (for $i=0$),
$\delta(q_1,\pmb{b},B)=\langle q_1,BB\rangle$, $\delta(q_1,\pmb{c},B)=\langle q_2,\lambda\rangle$,
$\delta(q_2,\pmb{c},B)=\langle q_2,\lambda\rangle$, $\delta(q_2,\pmb{d},A)=\langle q_3,\lambda\rangle$,
$\delta(q_3,\pmb{d},A)=\langle q_3,\lambda\rangle$, $\delta(q_2,\lambda,Z)=\langle q_f,Z\rangle$, $\delta(q_3,\lambda,Z)=\langle q_f,Z\rangle$.

\item \textbf{PDA via empty stack}: Pop $Z$ as well; add $\delta(\cdot,\lambda,Z)=\langle q_f,\lambda\rangle$ with appropriate state.

\item \textbf{DPDA}: Yes. The language has a fixed block structure: $\pmb{a}^i,\pmb{b}^j,\pmb{c}^j,\pmb{d}^i$, so the DPDA deterministically processes each block. The stack remembers $i$ (via $A$s) and $j$ (via $B$s overlaid on top). This is deterministic.

\item \textbf{Counting automaton}: No; both $i$ and $j$ are unbounded.

\item \textbf{Grammar from PDA} (Definition 10.7): Applying Definition 10.7 to part (b) and simplifying yields the equivalent grammar
\[
S\rightarrow \pmb{a}S\pmb{d}\mid T,\qquad
T\rightarrow \pmb{b}T\pmb{c}\mid \lambda.
\]
\end{enumerate}

%%------------------------------------------------------------
\item \textbf{Find $G_{P_3}$ from Example 10.5 using Definition 10.7.}

Since Definition 10.7 applies to PDAs that accept by empty stack, first convert $P_3$ to an equivalent empty-stack PDA by Theorem 10.5. Applying Definition 10.7 to that machine and deleting useless variables yields an equivalent grammar for the language of Example 10.5:
\[
Z\rightarrow X\mid Y,
\]
\[
X\rightarrow \pmb{a}X\pmb{b}\mid \pmb{ab},
\qquad
Y\rightarrow \pmb{aa}Y\pmb{b}\mid \pmb{aab}.
\]
Here $X$ generates $\{\pmb{a}^n\pmb{b}^n\mid n\geq 1\}$ and $Y$ generates $\{\pmb{a}^{2m}\pmb{b}^m\mid m\geq 1\}$, so $Z$ generates exactly the language recognized by $P_3$.

%%------------------------------------------------------------
\item \textbf{Prove by induction (Theorem 10.3): $A^{st}\DRightarrow x \Leftrightarrow \langle s,x,A\rangle\Dvdash\langle t,\lambda,\lambda\rangle$.}

\textbf{Claim} ($P(n)$): For all $n$, for all states $s,t$, stack symbol $A$, string $x$: $A^{st}\DRightarrow x$ in $n$ steps iff $\langle s,x,A\rangle\Dvdash\langle t,\lambda,\lambda\rangle$ in $n$ moves.

\textbf{Base} ($n=1$): $A^{st} \to a$ (terminal production from grammar of Definition 10.7) corresponds to $\delta(s,a,A) \ni \langle t,\lambda\rangle$, giving $\langle s,a,A\rangle\vdash\langle t,\lambda,\lambda\rangle$. \checkmark

\textbf{Step}: Assume $P(k)$ for all $k < n$. A derivation $A^{st}\DRightarrow x$ of length $n$ begins with $A^{st} \Rightarrow a B_1^{s_0 s_1} B_2^{s_1 s_2}\cdots B_m^{s_{m-1} t}$ (using a production from $\delta(s,a,A)\ni\langle s_0,B_1\cdots B_m\rangle$). Then each $B_i^{s_{i-1}s_i}$ independently derives $x_i$ in fewer than $n$ steps. By IH, $\langle s_{i-1},x_i,B_i\rangle\Dvdash\langle s_i,\lambda,\lambda\rangle$. Chaining: $\langle s,ax_1\cdots x_m,A\rangle\vdash\langle s_0,x_1\cdots x_m,B_1\cdots B_m\rangle\Dvdash\langle t,\lambda,\lambda\rangle$. So $x = ax_1\cdots x_m$ and $\langle s,x,A\rangle\Dvdash\langle t,\lambda,\lambda\rangle$. \qed

%%------------------------------------------------------------
\item \textbf{$L=\{\pmb{a}^n\pmb{b}^n\pmb{c}^m\pmb{d}^m\}\cup\{\pmb{a}^i\pmb{b}^j\pmb{c}^j\pmb{d}^i\}$.}

\begin{enumerate}
\item \textbf{PDA via final state}: Use nondeterminism to choose which language the input belongs to.
$P$: from initial state, nondeterministically branch to handle either $\pmb{a}^n\pmb{b}^n\pmb{c}^m\pmb{d}^m$ (Exercise 10.8) or $\pmb{a}^i\pmb{b}^j\pmb{c}^j\pmb{d}^i$ (Exercise 10.11).

\item \textbf{PDA via empty stack}: Start with a new initial state and make a $\lambda$-choice to the empty-stack PDA from Exercise 10.8 or the one from Exercise 10.11. This is the standard PDA construction for union.

\item \textbf{DPDA}: No. This is the classic ambiguous union
\[
\{\pmb{a}^n\pmb{b}^n\pmb{c}^m\pmb{d}^m\mid n,m\geq 0\}\cup
\{\pmb{a}^i\pmb{b}^j\pmb{c}^j\pmb{d}^i\mid i,j\geq 0\}.
\]
Every word $\pmb{a}^n\pmb{b}^n\pmb{c}^n\pmb{d}^n$ belongs to both halves, so the language has two incompatible structural readings on an infinite family of words. Since deterministic context-free languages are unambiguous, this language cannot be a DCFL.

\item \textbf{Counting automaton}: No; the string lengths are unbounded.

\item \textbf{Grammar from PDA}: An equivalent grammar is
\[
S\rightarrow AB\mid C,
\]
\[
A\rightarrow \pmb{a}A\pmb{b}\mid \lambda,\qquad
B\rightarrow \pmb{c}B\pmb{d}\mid \lambda,
\]
\[
C\rightarrow \pmb{a}C\pmb{d}\mid T,\qquad
T\rightarrow \pmb{b}T\pmb{c}\mid \lambda.
\]
\end{enumerate}

%%------------------------------------------------------------
\item \textbf{Find $G_{P_G}$ from Example 10.6 using Definition 10.7.}

Since $P_G$ has only one state $s$, the nonterminals are $Z$, $S^{ss}$, $\pmb{a}^{ss}$, and $\pmb{b}^{ss}$. Definition 10.7 gives
\[
Z\rightarrow S^{ss},
\]
\[
S^{ss}\rightarrow \pmb{a}S^{ss}\pmb{b}^{ss}\mid \pmb{a}\pmb{b}^{ss},
\]
\[
\pmb{b}^{ss}\rightarrow \pmb{b}.
\]
The variable $\pmb{a}^{ss}$ is useless and may be deleted. Renaming $S^{ss}$ as $S$ and $\pmb{b}^{ss}$ as $\pmb{b}$ gives the familiar grammar
\[
S\rightarrow \pmb{a}S\pmb{b}\mid \pmb{ab}.
\]

%%------------------------------------------------------------
\item \textbf{Prove (Theorem 10.4): $\langle s,xy,\alpha\rangle\Dvdash\langle s,y,\beta\rangle$ in $P$ iff $\langle s,xy,\alpha Y\rangle\Dvdash\langle s,y,\beta Y\rangle$ in $P_f$.}

\textbf{Proof by induction on the number of moves}:

\textbf{Base} (0 moves): Both sides are trivially equivalent (same configuration).

\textbf{Step}: If $\langle s,xy,\alpha\rangle\vdash\langle s',x'y,\alpha'\rangle\Dvdash\langle s,y,\beta\rangle$ in $P$:
The first move in $P$ corresponds to $\delta_P(s,a,A)\ni\langle s',\gamma\rangle$ (reading $a$ or $\lambda$, popping $A$, pushing $\gamma$). In $P_f$, the same move is available: $\delta_{P_f}$ includes all moves of $\delta_P$ plus new moves for the marker $Y$. The marker $Y$ sits below $\alpha$ and is not affected by moves on $\alpha$. By IH on the remaining $k-1$ moves, the result follows for $P_f$ with $Y$ appended. \qed

%%------------------------------------------------------------
\item \textbf{Prove (Theorem 10.5): same correspondence for $P_\lambda$.}

\textbf{Proof by induction}: Symmetric to Exercise 10.16; the marker $Y$ at the bottom of the stack (representing ``no more stack to pop'') is unaffected by all moves above it, so the computation in $P$ with stack $\alpha$ corresponds exactly to the computation in $P_\lambda$ with stack $\alpha Y$. \qed

%%------------------------------------------------------------
\item \textbf{$\{x\in\{\pmb{a},\pmb{b},\pmb{c}\}^* \mid |x|_{\pmb{a}}=|x|_{\pmb{b}}\wedge|x|_{\pmb{b}}>|x|_{\pmb{c}}\}$ is not context free.}

\textbf{Proof using closure properties}:

Let $K = \{x \mid |x|_{\pmb{a}}=|x|_{\pmb{b}}\wedge|x|_{\pmb{b}}>|x|_{\pmb{c}}\}$.

Suppose $K \in \PSIG$. Consider the regular language $R = \pmb{a}^*\pmb{b}^*\pmb{c}^*$ (strings of the form $\pmb{a}^i\pmb{b}^j\pmb{c}^k$).

$K \cap R = \{\pmb{a}^n\pmb{b}^n\pmb{c}^k \mid n>k\}$.

If $K$ were a CFL, then $K \cap R$ would be a CFL (by Theorem 10.6: CFL $\cap$ regular $=$ CFL). But $\{\pmb{a}^n\pmb{b}^n\pmb{c}^k \mid n>k\}$ is not context free (one can apply the pumping lemma to strings $\pmb{a}^p\pmb{b}^p\pmb{c}^{p-1}$; any pumping either changes the $\pmb{a}$-count differently from the $\pmb{b}$-count, or changes the $\pmb{c}$-count, in each case violating the constraints).

So $K$ is not a CFL. \qed

%%------------------------------------------------------------
\item \textbf{Two-tape automaton.}

\begin{enumerate}
\item \textbf{Transition function}: For a two-tape automaton with tapes over $\Sigma$, the state transition function $\delta$ has domain $S \times (\Sigma\cup\{\pmb{B}\}) \times (\Sigma\cup\{\pmb{B}\})$ (current state, symbol read from tape 1, symbol read from tape 2) and range $S \times (\Sigma\cup\{\pmb{B}\}) \times \{L,R,S\} \times (\Sigma\cup\{\pmb{B}\}) \times \{L,R,S\}$ (new state, write to tape 1, move head 1, write to tape 2, move head 2).

\item \textbf{Two-tape automaton for $\{\pmb{a}^n\pmb{b}^n\pmb{c}^n \mid n\geq 1\}$}: Let
\[
M=\langle\{\pmb{a},\pmb{b},\pmb{c}\},\Gamma,\{q_0,q_1,q_2,q_3,q_4,q_f,q_r\},q_0,\delta,q_f\rangle,
\]
where tape 1 initially holds the input and tape 2 is blank, and $\Gamma=\{\pmb{a},\pmb{b},\pmb{c},X,\pmb{B}\}$. The intended behavior is:
\begin{enumerate}[1.]
\item In $q_0$, scan the $\pmb{a}$-block on tape 1. For each $\pmb{a}$ read on tape 1, write an $X$ on tape 2 and move both heads right. If no $\pmb{a}$ is seen, reject. On the first $\pmb{b}$, leave tape 1 where it is, move tape 2 left into state $q_1$, and rewind tape 2 to its leftmost $X$.
\item In $q_2$, compare the $\pmb{b}$-block with the stored $X$'s: on $(\pmb{b},X)$ move right on both tapes. If the tape-1 symbol changes to $\pmb{c}$ exactly when tape 2 reaches blank, go to $q_3$; otherwise reject.
\item In $q_3$, hold tape 1 on the first $\pmb{c}$ and rewind tape 2 left to its leftmost $X$ again.
\item In $q_4$, compare the $\pmb{c}$-block with the same $X$'s: on $(\pmb{c},X)$ move right on both tapes. Accept exactly when both heads simultaneously reach blank; any other mismatch or extra symbol causes rejection.
\end{enumerate}
Thus tape 2 stores the number $n$ once, and the machine checks both the $\pmb{b}$-block and the $\pmb{c}$-block against that same count.
\end{enumerate}

%%------------------------------------------------------------
\item \textbf{$\{\pmb{a}^n\pmb{b}^n\pmb{c}^n\}$ is not context free; $\{x \mid |x|_{\pmb{a}}=|x|_{\pmb{b}}\}$ is not context free.}

\begin{enumerate}
\item \textbf{$\{\pmb{a}^n\pmb{b}^n\pmb{c}^n\}$ is not CFL}: Assume it is, with pumping constant $N$. Take $s=\pmb{a}^N\pmb{b}^N\pmb{c}^N$. By the pumping lemma, $s=uvwxy$ with $|vwx|\leq N$ and $|vx|\geq 1$.

Since $|vwx|\leq N$, the substring $vwx$ cannot span all three blocks. So $v$ and $x$ together involve at most two of $\{\pmb{a},\pmb{b},\pmb{c}\}$. Pumping changes the count of at most two symbols while leaving the third unchanged. But all three must be equal, so the pumped string fails $|u|_{\pmb{a}}=|u|_{\pmb{b}}=|u|_{\pmb{c}}$ for some $i\neq 1$. Contradiction. \qed

\item \textbf{Disproof}: the printed statement is false, because
\[
L = \{x \mid |x|_{\pmb{a}}=|x|_{\pmb{b}}\}
\]
is context free. For example, the grammar
\[
S \;\to\; \pmb{c}S \;\mid\; SS \;\mid\; \pmb{a}S\pmb{b} \;\mid\; \pmb{b}S\pmb{a} \;\mid\; \lambda
\]
generates exactly the strings over $\{\pmb{a},\pmb{b},\pmb{c}\}$ having equally many $\pmb{a}$'s and $\pmb{b}$'s. Therefore the claim to be proved is false.
\end{enumerate}

%%------------------------------------------------------------
\item \textbf{DPDA constructions and non-DCFL example.}

\begin{enumerate}
\item DPDA for $\{\pmb{c}^n\pmb{b}^m \mid n\geq 1, n=m\} \cup \{\pmb{a}^n\pmb{b}^m \mid n\geq 1, n=2m\}$:

Let
\[
P=\langle\{\pmb{a},\pmb{b},\pmb{c}\},\{Z,A\},
\{q_0,q_c,q_{a,1},q_{a,0},q_b,q_f\},q_0,\delta,Z,\{q_f\}\rangle.
\]
The first symbol chooses the branch deterministically.

For the $\pmb{c}^n\pmb{b}^n$ branch, use the standard one-counter routine:
\begin{align*}
\delta(q_0,\pmb{c},Z)&=\langle q_c,AZ\rangle,\\
\delta(q_c,\pmb{c},A)&=\langle q_c,AA\rangle,\\
\delta(q_c,\pmb{b},A)&=\langle q_b,\lambda\rangle,\\
\delta(q_b,\pmb{b},A)&=\langle q_b,\lambda\rangle,\\
\delta(q_b,\lambda,Z)&=\langle q_f,Z\rangle.
\end{align*}

For the $\pmb{a}^n\pmb{b}^m$ branch with $n=2m$, the machine records the parity of the number of $\pmb{a}$'s: state $q_{a,1}$ means an odd number of $\pmb{a}$'s has been seen since the last push, and state $q_{a,0}$ means an even number.
\begin{align*}
\delta(q_0,\pmb{a},Z)&=\langle q_{a,1},Z\rangle,\\
\delta(q_{a,1},\pmb{a},Z)&=\langle q_{a,0},AZ\rangle,\\
\delta(q_{a,1},\pmb{a},A)&=\langle q_{a,0},AA\rangle,\\
\delta(q_{a,0},\pmb{a},A)&=\langle q_{a,1},A\rangle,\\
\delta(q_{a,0},\pmb{b},A)&=\langle q_b,\lambda\rangle.
\end{align*}
Every second $\pmb{a}$ pushes one $A$, so after reading $\pmb{a}^n$ the stack contains exactly $n/2$ copies of $A$. The $\pmb{b}$-phase then pops one $A$ per $\pmb{b}$. Acceptance occurs only when the input ends and the stack returns to $Z$, so this branch accepts exactly $\{\pmb{a}^{2m}\pmb{b}^m\mid m\geq 1\}$.

\item Define $h: \{\pmb{a},\pmb{b},\pmb{c}\}^* \to \{\pmb{a},\pmb{b}\}^*$ by $h(\pmb{a})=\pmb{a}$, $h(\pmb{b})=\pmb{b}$, $h(\pmb{c})=\pmb{a}$. Then:
$h(\{\pmb{c}^n\pmb{b}^m\mid n=m\}) = \{\pmb{a}^n\pmb{b}^n\mid n\geq 1\}$ and $h(\{\pmb{a}^n\pmb{b}^m\mid n=2m\}) = \{\pmb{a}^n\pmb{b}^m\mid n=2m\}$.
$h(L) = \{\pmb{a}^n\pmb{b}^n\}\cup\{\pmb{a}^n\pmb{b}^m\mid n=2m\}$. The discussion immediately before Theorem 10.9 identifies this union as a context-free language that is \emph{not} deterministic context free. Hence this homomorphism takes a DCFL to a non-DCFL, as required.
\end{enumerate}

%%------------------------------------------------------------
\item \textbf{Ogden's lemma: $\{\pmb{a}^k\pmb{b}^n\pmb{c}^m \mid k\neq n \wedge n\neq m\}$ is not CFL.}

Assume $L = \{\pmb{a}^k\pmb{b}^n\pmb{c}^m \mid k\neq n, n\neq m\}$ is CFL with Ogden constant $N$.

Choose $p > N$. Let $s = \pmb{a}^{p!+p}\pmb{b}^p\pmb{c}^{p!+p} \in L$ (since $p!+p \neq p$). Mark all $\pmb{b}$-symbols.

By Ogden's lemma: $s = uvwxy$ with $vwx$ containing at most $N$ marked positions, $vx$ has at least one marked position, all pumped strings in $L$.

Since all marks are on $\pmb{b}$s, and $|vwx|\leq N \leq p$, $vwx$ lies within the $\pmb{b}$-block (possibly with some $\pmb{a}$s or $\pmb{c}$s on the boundary). Since $vx$ has at least one $\pmb{b}$, $|vx|_{\pmb{b}} \geq 1$; let $|vx|_{\pmb{b}} = j$ ($1\leq j\leq N$).

For the pumped string $uv^i wx^i y$: the count of $\pmb{b}$s is $p + (i-1)j$. The counts of $\pmb{a}$s and $\pmb{c}$s may also change if $v$ or $x$ extend into the $\pmb{a}$ or $\pmb{c}$ blocks.

\textbf{Case 1}: $vx$ contains only $\pmb{b}$s. Then $k = p!+p$ and $m = p!+p$ for all $i$. Choose $i$ such that $p + (i-1)j = p!+p$: then $(i-1)j = p!$, so $i-1 = p!/j$. Since $j\leq N\leq p-1$, $j\mid p!$, so $i$ is an integer. The pumped string has $|\pmb{b}|=p!+p=k=m$, violating $k\neq n$ or $n\neq m$. Contradiction.

\textbf{Case 2}: $vx$ spans the $\pmb{a}$/$\pmb{b}$ boundary. Let $|vx|_{\pmb{a}}=\ell \geq 0$ and $|vx|_{\pmb{b}}=j \geq 1$ (at least one marked $\pmb{b}$). Choose $i - 1 = p!/j$ (an integer since $j \leq N < p$). Then the pumped string has $n = p + p! = p!+p$ and $m = p!+p$ (the $\pmb{c}$-block is untouched), so $n = m$, violating $n\neq m$. Contradiction.

\textbf{Case 3}: $vx$ spans the $\pmb{b}$/$\pmb{c}$ boundary. Let $|vx|_{\pmb{b}}=j \geq 1$ and $|vx|_{\pmb{c}}=r \geq 0$. Choose $i-1 = p!/j$. The pumped string has $n = p!+p$ and $k = p!+p$ (the $\pmb{a}$-block is untouched), so $k = n$, violating $k\neq n$. Contradiction. \qed

%%------------------------------------------------------------
\item \textbf{Prove (Theorem 10.6) by induction.}

Theorem 10.6 states that if $P_1$ is a PDA and $A_2$ is a DFA (or regular language), then their ``intersection'' machine $P^{\cap}$ (Product machine) satisfies the correspondence stated.

\textbf{Proof by induction on the number of moves}:

\textbf{Base} (0 moves): $\langle\langle s_1,s_2\rangle,y,\alpha\rangle\Dvdash\langle\langle s_1,s_2\rangle,y,\alpha\rangle$ (trivially). Here $t_1=s_1$, $t_2=s_2$, $x=\lambda$, so $\DBAR_2(s_2,\lambda)=s_2=t_2$. \checkmark

\textbf{Step}: If the product machine takes $k+1$ moves from $\langle\langle s_1,s_2\rangle,xy,\alpha\rangle$: the first move reads some $a$ (or $\lambda$) and moves to $\langle\langle s_1',s_2'\rangle,\cdot,\cdot\rangle$. In the product machine, this corresponds to a joint transition: $P_1$ moves from $s_1$ and $A_2$ reads $a$ to move from $s_2$. By IH on $k$ moves, the remaining computation in the product machine corresponds to $P^{\cap}$ and $\DBAR_2$. The $A_2$ component is precisely $\DBAR_2(s_2,x)$ after reading $x$. \qed

%%------------------------------------------------------------
\item \textbf{Prove: DPDA $P$ has equivalent DPDA $P'$ with total transitions, no empty stack, scans full input.}

\begin{enumerate}
\item \textbf{Never empty stack}: First place a permanent bottom marker $\pmb{\$}$ under the original start symbol. Introduce a new start state $s_0$ with the single move $\delta(s_0,\lambda,\pmb{\$})=\langle s,Z\pmb{\$}\rangle$, where $s$ is the old start state and $Z$ the old initial stack symbol. Whenever the old machine would have popped its last stack symbol, the new machine instead reaches a designated state whose stack top is $\pmb{\$}$. Thus the stack is never empty.

\item \textbf{Always scans full input}: Add an endmarker $\pmb{\cent}$ to the input alphabet. Replace every old accepting halt by a move into an accepting scan state $q_{\mathrm{acc}}$, and every rejecting halt by a move into a rejecting scan state $q_{\mathrm{rej}}$. In both scan states, on every ordinary input symbol the machine simply moves right and leaves the stack unchanged; on $\pmb{\cent}$ it stops. Hence every computation reads the entire input before it halts.

\item \textbf{Total transitions}: Now fill in whatever cases remain undefined, but in a way that respects the DPDA rules. If for some pair $(q,X)$ there is already a $\lambda$-move, then no input move may be added there; the existing $\lambda$-move already supplies the unique next step. If there is no $\lambda$-move for $(q,X)$, then for each missing input symbol add a transition to $q_{\mathrm{rej}}$ that preserves the stack. If neither input moves nor a $\lambda$-move existed for $(q,X)$, add the single $\lambda$-move to $q_{\mathrm{rej}}$. Thus every reachable configuration has a next move, and no input/$\lambda$ conflict is introduced.

These three repairs preserve the accepted language and keep the machine deterministic, so the resulting DPDA has all the required properties. \qed
\end{enumerate}

%%------------------------------------------------------------
\item \textbf{\ASIG\ (DCFL) is closed under complementation.}

Using Exercise 10.24: For a DPDA $P$ (accepting via final state), $P$ can be made total (all transitions defined), never emptying stack, and always scanning full input. Under these conditions:

Every computation on every input string either ends in a final state (accept) or a nonfinal state (reject). There are no infinite loops or premature halts.

Construct $P'$ by swapping final and nonfinal states of $P$: a string is accepted by $P'$ iff it is rejected by $P$. The main issue (mentioned in the hint) is $\lambda$-moves cycling through both final and nonfinal states: since Exercise 10.24 ensures the DPDA is well-behaved (complete, non-emptying, full-input scanning), after processing any input string, the computation ends deterministically in exactly one state. Swapping final/nonfinal gives the complement. \qed

%%------------------------------------------------------------
\item \textbf{Example: \ASIG\ not closed under concatenation.}

Let
\[
L_1=\{\pmb{0}^i\pmb{1}^i\pmb{2}^k\mid i,k\geq 0\}
\qquad\text{and}\qquad
L_2=\{\pmb{0}^i\pmb{1}^k\pmb{2}^k\mid i,k\geq 0\}.
\]
Both are DCFLs, but $L_1\cup L_2$ is not. Now define
\[
L_3=\pmb{a}L_1\cup L_2.
\]
$L_3$ is a DCFL, because the initial $\pmb{a}$ tells the DPDA whether to use the $L_1$-routine or the $L_2$-routine. The regular language $\pmb{a}^*$ is certainly a DCFL. Suppose \ASIG\ were closed under concatenation. Then $\pmb{a}^*L_3$ would be a DCFL, and therefore so would
\[
L_4=\pmb{a}^*L_3\cap \pmb{a}\pmb{0}^*\pmb{1}^*\pmb{2}^*.
\]
But
\[
L_4=\pmb{a}L_1\cup \pmb{a}L_2.
\]
If a DPDA recognized $L_4$, we could recognize $L_1\cup L_2$ by simulating that DPDA after an imaginary initial $\pmb{a}$. Since $L_1\cup L_2$ is not a DCFL, $L_4$ cannot be a DCFL, and so \ASIG\ is not closed under concatenation.

%%------------------------------------------------------------
\item \textbf{Example: \ASIG\ not closed under Kleene closure.}

With $L_3$ as in Exercise 10.26, let
\[
L_5=\{\pmb{a}\}\cup L_3.
\]
$L_5$ is a DCFL. Suppose \ASIG\ were closed under Kleene closure. Then $L_5^*$ would be a DCFL, and hence so would
\[
L_5^*\cap \pmb{a}\pmb{0}^*\pmb{1}^*\pmb{2}^*.
\]
The only strings of $L_5^*$ that have the form $\pmb{a}\pmb{0}^*\pmb{1}^*\pmb{2}^*$ are
\[
\{\pmb{a}\}\cup \pmb{a}L_1\cup \pmb{a}L_2.
\]
Intersecting further with the regular complement of $\{\pmb{a}\}$ would leave $\pmb{a}L_1\cup \pmb{a}L_2$, which is not a DCFL by the proof of Exercise 10.26. Therefore $L_5^*$ cannot be a DCFL, and \ASIG\ is not closed under Kleene closure.

%%------------------------------------------------------------
\item \textbf{$\{\pmb{ca}^n\pmb{b}^m \mid n=m\}\cup\{\pmb{a}^n\pmb{b}^m \mid n=2m\}$ is a DCFL.}

A DPDA can distinguish the two cases by the first symbol. On the $\pmb{c}$-branch, it first consumes the initial $\pmb{c}$ and then uses the standard DPDA for $\{\pmb{a}^n\pmb{b}^n\}$: push one $A$ per $\pmb{a}$ and pop one per $\pmb{b}$. On the $\pmb{a}$-branch, it uses the parity-counting construction from Exercise 10.20(a): every second $\pmb{a}$ pushes one $A$, and each $\pmb{b}$ pops one $A$. Thus the second branch accepts exactly $\{\pmb{a}^{2m}\pmb{b}^m\mid m\geq 1\}$. Since the first symbol determines the branch uniquely, the whole machine is deterministic. \qed

%%------------------------------------------------------------
\item \textbf{$L_1 - R_2$ is context free if $L_1$ is CFL and $R_2$ is regular.}

\begin{enumerate}
\item \textbf{Via product machine}: Extend the intersection proof (Theorem 10.6). Construct a product machine where the PDA simulates $L_1$ and the DFA simulates $R_2$; modify final states to accept when the PDA is in a final state and the DFA is in a \emph{nonfinal} state. This gives $L_1 \cap \overline{R_2} = L_1 - R_2$.

\item \textbf{Via closure}: $L_1 - R_2 = L_1 \cap \overline{R_2}$. Since regular languages are closed under complement, $\overline{R_2}$ is regular. Since CFLs are closed under intersection with regular languages (Theorem 10.6), $L_1 \cap \overline{R_2}$ is a CFL. \qed
\end{enumerate}

%%------------------------------------------------------------
\item \textbf{$L_1\cup R_2$ is context free if $L_1$ is CFL and $R_2$ is regular.}

\begin{enumerate}
\item \textbf{Via modified product machine}: Accept when the PDA is in a final state (for $L_1$) OR the DFA is in a final state (for $R_2$). This is a product machine with final states $F_1 \times S_2 \cup S_1 \times F_2$.

\item \textbf{Via closure}: $L_1 \cup R_2$: since $R_2$ is regular, it is also a CFL. CFLs are closed under union: $L_1 \cup R_2 \in \PSIG$. \qed
\end{enumerate}

%%------------------------------------------------------------
\item \textbf{If $L_1$ is DCFL and $R_2$ is regular, $R_2 - L_1$ is DCFL.}

$R_2 - L_1 = R_2 \cap \overline{L_1}$. Since $L_1$ is DCFL, $\overline{L_1}$ is DCFL (Exercise 10.25). Since $R_2$ is regular, and the discussion immediately after Theorem 10.6 notes that the product construction remains deterministic when the PDA factor is deterministic, $R_2 \cap \overline{L_1}$ is a DCFL. \qed

%%------------------------------------------------------------
\item \textbf{$\{w\pmb{2}w^r\}$ is DCFL; $\{ww^r\}$ is not DCFL.}

\begin{enumerate}
\item \textbf{$\{w\pmb{2}w^r\}$ is DCFL}: The DPDA pushes all symbols before $\pmb{2}$ onto the stack, then pops matching symbols for each character after $\pmb{2}$ (reversed). Since $\pmb{2}$ marks the center deterministically, the DPDA knows exactly when to switch from push to pop mode. \qed

\item \textbf{$\{ww^r\}$ is not DCFL}: Let
\[
E=\{ww^r\mid w\in\{\pmb{0},\pmb{1}\}^*\}.
\]
Assume $E$ were a DCFL. By Exercise 10.13, DCFLs are closed under inverse homomorphism. Define a homomorphism
\[
h(\pmb{a})=\pmb{01},\qquad h(\pmb{b})=\pmb{1010},\qquad h(\pmb{c})=\pmb{0101},
\]
and let $R=\pmb{c}^+\pmb{b}^+\cup \pmb{a}^+\pmb{b}^+$. For $x=\pmb{c}^n\pmb{b}^m$ we have
\[
h(x)=(\pmb{01})^{2n}(\pmb{10})^{2m},
\]
which is an even palindrome iff $n=m$. For $x=\pmb{a}^n\pmb{b}^m$ we have
\[
h(x)=(\pmb{01})^n(\pmb{10})^{2m},
\]
which is an even palindrome iff $n=2m$. Therefore
\[
R\cap h^{-1}(E)=
\{\pmb{c}^n\pmb{b}^n\mid n\geq 1\}\cup
\{\pmb{a}^n\pmb{b}^m\mid n\geq 1,\ n=2m\}.
\]
But Exercise 10.20 and the discussion immediately before Theorem 10.9 identify this language as not deterministic context free. This contradiction shows that $E$ is not a DCFL. \qed
\end{enumerate}

%%------------------------------------------------------------
\item \textbf{Non-closure of \ASIG\ under various operations.}

\begin{enumerate}
\item $L_1\cdot L_2$ need not be DCFL: Example from Exercise 10.26.
\item $L_1 - L_2$ need not be DCFL: Let
\[
L_1=\{\pmb{a}^n\pmb{b}^n\pmb{c}^m\mid n,m\geq 0\}
\qquad\text{and}\qquad
K=\{\pmb{a}^m\pmb{b}^n\pmb{c}^n\mid m,n\geq 0\}.
\]
Both $L_1$ and $K$ are DCFLs. By Exercise 10.25, $\sim\!K$ is also a DCFL. Now
\[
L_1-(\sim\!K)=L_1\cap K=\{\pmb{a}^n\pmb{b}^n\pmb{c}^n\mid n\geq 0\},
\]
which is not even context free. So difference does not preserve \ASIG.
\item $L_1^*$ need not be DCFL: Exercise 10.27.
\item $L_1^r$ need not be DCFL: let
\[
L=\{\pmb{a}\pmb{2}^j\pmb{1}^i\pmb{0}^i\mid i,j\geq 0\}\cup
\{\pmb{b}\pmb{2}^j\pmb{1}^j\pmb{0}^i\mid i,j\geq 0\}.
\]
This is a DCFL, because the first symbol determines which deterministic routine to use. But
\[
L^r=\{\pmb{0}^i\pmb{1}^i\pmb{2}^j\pmb{a}\mid i,j\geq 0\}\cup
\{\pmb{0}^i\pmb{1}^j\pmb{2}^j\pmb{b}\mid i,j\geq 0\}.
\]
If $L^r$ were a DCFL, then the same one-letter-lookahead argument as in Exercise 10.30 would turn it into a DPDA for
\[
\{\pmb{0}^i\pmb{1}^i\pmb{2}^j\mid i,j\geq 0\}\cup
\{\pmb{0}^i\pmb{1}^j\pmb{2}^j\mid i,j\geq 0\},
\]
which is the same standard non-DCFL union pattern discussed before Theorem 10.9, now with an extra trailing block. So reversal does not preserve \ASIG.
\end{enumerate}

%%------------------------------------------------------------
\item \textbf{Closure of \PSIG\ and \ASIG\ under quotient.}

\begin{enumerate}
\item \textbf{\PSIG\ is not closed under quotient}: Let
\[
L=\{\pmb{a}^n\pmb{b}^n\pmb{c}^{n+m}\pmb{d}^m\mid n,m\geq 0\}
\qquad\text{and}\qquad
K=\{\pmb{c}^m\pmb{d}^m\mid m\geq 0\}.
\]
The language $K$ is context free. So is $L$: a PDA pushes one marker for each $\pmb{a}$, and on each $\pmb{b}$ replaces one such marker by a marker that must later be matched by a $\pmb{c}$. After those forced $\pmb{c}$'s have been read, any additional $\pmb{c}$'s are counted against the final $\pmb{d}$'s. Now
\[
L/K=\{\pmb{a}^n\pmb{b}^n\pmb{c}^n\mid n\geq 0\}.
\]
Indeed, every $\pmb{a}^n\pmb{b}^n\pmb{c}^n$ is in the quotient by taking the suffix $\lambda\in K$. Conversely, if $x\in L/K$, then $xy\in L$ for some $y=\pmb{c}^m\pmb{d}^m\in K$. Since every word of $L$ has the form $\pmb{a}^n\pmb{b}^n\pmb{c}^{n+r}\pmb{d}^r$, the suffix $\pmb{c}^m\pmb{d}^m$ can only begin after exactly the first $n$ $\pmb{c}$'s; otherwise the number of $\pmb{c}$'s preceding the $\pmb{d}$-block would exceed the number of trailing $\pmb{d}$'s. Hence $x=\pmb{a}^n\pmb{b}^n\pmb{c}^n$. But $\{\pmb{a}^n\pmb{b}^n\pmb{c}^n\mid n\geq 0\}$ is not context free, so \PSIG\ is not closed under quotient.

\item \textbf{\ASIG\ is not closed under quotient}: The same counterexample works for deterministic context-free languages, because both $L$ and $K$ are accepted by DPDAs. For $K$, use the standard DPDA for $\{\pmb{c}^m\pmb{d}^m\}$. For $L$, use the deterministic phase structure described above:
\[
\pmb{a}\text{-phase} \rightarrow \pmb{b}\text{-phase} \rightarrow
\text{forced }\pmb{c}\text{-phase} \rightarrow \text{extra }\pmb{c}\text{-phase}
\rightarrow \pmb{d}\text{-phase}.
\]
Thus $L,K\in\ASIG$, but
\[
L/K=\{\pmb{a}^n\pmb{b}^n\pmb{c}^n\mid n\geq 0\}
\]
is not even context free. Therefore \ASIG\ is not closed under quotient.
\end{enumerate}

%%------------------------------------------------------------
\item \textbf{Closure under operator $b$ (Theorem 5.11).}

\begin{enumerate}
\item \PSIG\ is \textbf{not} closed under $b$. Let
\[
A=\{\pmb{a}^m\pmb{b}^m\pmb{c}^n\pmb{\#\#}\pmb{d}^{3n}\mid m,n\geq 0\}.
\]
$A$ is context free; for example,
\[
S\rightarrow TU,\qquad
T\rightarrow \pmb{a}T\pmb{b}\mid \lambda,\qquad
U\rightarrow \pmb{c}U\pmb{ddd}\mid \pmb{\#\#}.
\]
Now intersect $A^b$ with the regular language $\pmb{a}^*\pmb{b}^*\pmb{c}^*\pmb{\#}$.
If $x=\pmb{a}^i\pmb{b}^j\pmb{c}^k\pmb{\#}$ belongs to this intersection, then some $y$
of the same length satisfies $xy\in A$. The only possible choice is
\[
xy=\pmb{a}^n\pmb{b}^n\pmb{c}^n\pmb{\#\#}\pmb{d}^{3n}
\]
for some $n$, so $i=j=k=n$. Conversely, every $\pmb{a}^n\pmb{b}^n\pmb{c}^n\pmb{\#}$
does arise this way. Hence
\[
A^b\cap \pmb{a}^*\pmb{b}^*\pmb{c}^*\pmb{\#}
=\{\pmb{a}^n\pmb{b}^n\pmb{c}^n\pmb{\#}\mid n\geq 0\},
\]
which is not context free. Therefore \PSIG\ is not closed under $b$.

\item \ASIG\ is \textbf{not} closed under $b$. The same language $A$ is accepted by a DPDA: deterministically match the $\pmb{a}$'s against the $\pmb{b}$'s, then push three markers for each $\pmb{c}$, read the fixed delimiter $\pmb{\#\#}$, and pop one marker for each $\pmb{d}$. Thus $A\in\ASIG$, but $A^b$ is not even context free by part (a). So \ASIG\ is not closed under $b$.
\end{enumerate}

%%------------------------------------------------------------
\item \textbf{Closure under operator $Y$ (Theorem 5.7).}

\begin{enumerate}
\item \PSIG\ is closed under $Y$. If $P$ is a PDA for $L$, build a new PDA with two copies of each state. In the first copy it simulates $P$ normally, but once during the computation it may read one input letter and move to the second copy without changing the simulated stack. That single move represents the deleted letter. After that point the second copy continues to simulate $P$ exactly. The new PDA therefore accepts precisely those strings obtained by deleting exactly one letter from a string of $L$.

\item \ASIG\ is \textbf{not} closed under $Y$. Let
\[
L=\{w\pmb{2}w^r\mid w\in\{\pmb{0},\pmb{1}\}^*\}.
\]
By Exercise 10.30(a), $L$ is a DCFL. If we delete exactly one symbol from a word of $L$, the resulting string has no $\pmb{2}$ exactly when the deleted symbol was the middle $\pmb{2}$. Hence
\[
Y(L)\cap\{\pmb{0},\pmb{1}\}^*=\{ww^r\mid w\in\{\pmb{0},\pmb{1}\}^*\}.
\]
By Exercise 10.30(b), the language on the right is not a DCFL. Since \ASIG\ is closed under intersection with regular sets, $Y(L)$ cannot be a DCFL. So \ASIG\ is not closed under $Y$.
\end{enumerate}

%%------------------------------------------------------------
\item \textbf{Closure under prefix operator $P$ (Exercise 5.16).}

\begin{enumerate}
\item \textbf{\PSIG\ is closed under $P$}: Let $M$ be a PDA for $L$. Build a new PDA $M'$ with two copies of each state of $M$. In the first copy, $M'$ simulates $M$ on the actual input. At any moment, $M'$ may take a $\lambda$-move to the second copy. In that second copy, every transition of $M$ that originally read a symbol of $\Sigma$ is replaced by a $\lambda$-transition with the same stack effect, and the original $\lambda$-moves are kept as $\lambda$-moves. Thus, once the real input has been exhausted, $M'$ may nondeterministically \emph{guess} a continuation and simulate $M$ on that guessed suffix. Therefore $M'$ accepts exactly those $x$ for which some $y$ satisfies $xy\in L$, that is, exactly $P(L)$.

\item \textbf{\ASIG\ is closed under $P$}: Let $P$ be a DPDA normalized as in Exercise 10.24. For each stack symbol $X$, use two tagged copies $X^{\mathrm{live}}$ and $X^{\mathrm{dead}}$. While simulating $P$, the new machine maintains on the stack not just the current stack symbols but also, for each stack position, whether the subconfiguration beginning at that position can still be extended to an accepting computation. Because the original machine is deterministic, this tag is forced by the current state and the tagged symbols that replace $X$ in the next move. The machine therefore updates the tags deterministically whenever it pushes or pops. After the input ends, the simulated configuration is live exactly when the topmost tag is $\mathrm{live}$, and that is exactly the condition $x\in P(L)$. Hence a DPDA can accept the prefix language.
\end{enumerate}

%%------------------------------------------------------------
\item \textbf{Closure under operator $F$ (Exercise 5.19).}

\begin{enumerate}
\item \textbf{\PSIG\ is not closed under $F$}. Let
\[
K=\{\pmb{a}^i\pmb{b}^j\pmb{c}^k\mid i=j\}\cup\{\pmb{a}^i\pmb{b}^j\pmb{c}^k\mid j=k\}.
\]
$K$ is context free, but by Exercise 10.23 its complement over $\{\pmb{a},\pmb{b},\pmb{c}\}^*$ is not context free. Introduce a new symbol $\pmb{\$}$ and let
\[
B=\{x\pmb{\$}\mid x\in\{\pmb{a},\pmb{b},\pmb{c}\}^*\}\cup\{x\pmb{\$\$}\mid x\in K\}.
\]
$B$ is context free. A word $x\pmb{\$}$ is in $F(B)$ exactly when $x\notin K$, because it fails to be maximal precisely when it can be extended to $x\pmb{\$\$}$. Therefore
\[
F(B)\cap\{\pmb{a},\pmb{b},\pmb{c}\}^*\pmb{\$}
=\{x\pmb{\$}\mid x\notin K\}.
\]
If $F(B)$ were context free, then intersecting with the regular language $\{\pmb{a},\pmb{b},\pmb{c}\}^*\pmb{\$}$ and deleting the final $\pmb{\$}$ would give $\sim\!K$ as a CFL, contradiction.

\item \textbf{\ASIG\ is closed under $F$}: Start with a DPDA normalized as in Exercise 10.24. For each input word $x$, let $C_x$ be the unique configuration reached after all of $x$ has been scanned. Then $x\in F(L)$ exactly when two conditions hold: $C_x$ is accepting, and $C_x$ is not \emph{prolific}, meaning that no nonempty continuation from $C_x$ can still lead to acceptance. This is the same configuration-liveness idea used in part (b) for the prefix operator, with ``live'' replaced by ``accepting but not prolific.'' A DPDA can therefore simulate the original machine and accept exactly the words in $F(L)$, so \ASIG\ is closed under $F$.
\end{enumerate}

\end{enumerate}
